\documentclass[a4paper]{article}
%\usepackage{simplemargins}

%\usepackage[square]{natbib}
\usepackage{amsmath}
\usepackage{amsfonts}
\usepackage{amssymb}
\usepackage{graphicx}
\usepackage{hyperref}
\usepackage{multirow}
\usepackage{makecell}

\begin{document}

\Large
 \begin{center}  
 Dynamic Risk Posture Evaluation Using Automated Threat Intelligence
\vspace{2em}
%\hspace{10pt}

% Author names and affiliations
\large


\end{center}

%\hspace{10pt}

\normalsize

One of the most urgent issues of contemporary enterprises is cybersecurity risk management, where manual, episodic assessments fail to consolidate real-time global threat intelligence with fine-grained technical telemetry. To address this, we describe DRPE (Dynamic Risk Posture Evaluation), a smart, fully containerized cybersecurity orchestration framework built using a high-performance FastAPI back-end, a React 18 enterprise dashboard, and a PostgreSQL relational database, coordinated by Docker Compose to be deployed in a single command.

Technically at the heart of the platform, DRPE uses a proprietary multi-variable Risk Engine which combines local technical telemetry of the OpenVAS/GVM infrastructure with global threat intelligence feeds of AbuseIPDB and AlienVault OTX. This generates a normalized, context-specific security posture score reflecting both intrinsic severity and real-world exploitability. Furthermore, the platform uses a Dynamic AI Engine with Google Gemini 1.5 Flash to support strategic decision-making, converting complex CVE identifiers and CVSS vectors to natural-language mitigation policies and executive-friendly impact assessments.

The system security is based on a stringent JWT authentication module alongside an immutable Activity Logging framework that provides auditability and proof of compliance. Client-side reporting (js PDF) generates executive compliance reports with negligible server load. Our experimental analysis indicates that the interactive 2D/3D Force-Directed Graph visualizations in the system significantly improve the speed and effectiveness of risk propagation analysis through network nodes as opposed to tabular lists of vulnerabilities.\\
%\hspace{10pt}


\vspace{1em}

\begin{center}
\begin{tabular}{ll}
{\textbf{Keywords:}} & Cybersecurity, DRPE, React, FastAPI, PostgreSQL  \\ \vspace{0.5em}
{\textbf{Project ID:}} & IT-F25-27             \\ \vspace{0.5em}
{\textbf{Project Area/ Domains:}} & Information Technology, Cybersecurity              \\ \vspace{0.5em} 
{\textbf{Supervisor Name:}} & Dr. Faisal Bashir Hussain    \\ \vspace{0.5em}
{\textbf{Supervisor Designation:}} & Dean, Faculty of Engineering Sciences\\ \vspace{0.5em}
{\textbf{Department:}} & Department of Information Technology    \\ \vspace{0.5em} 
{\textbf{Supervisor Email:}} & fbhussain@bahria.edu.pk\\
\end{tabular}
\end{center}

\vspace{1em}
\noindent\textbf{Group Members Detail:}\\

\begin{center}
\begin{tabular}{|p{0.4\textwidth}|p{0.4\textwidth}|} 
\hline
\textbf{Member 1} & \textbf{Member 2}  \\
\hline
\makecell{Zainab Khawar \\ +92-332-1569733\\ \href{https://pk.linkedin.com/in/zainab-khawar-a963aa30b}{Linkedin}}
&
\makecell{Adina Saif Ullah \\ +92-331-5113307\\ \href{https://www.linkedin.com/in/adina-saif-ullah-629863257}{Linkedin}}
  \\ 
\hline

\end{tabular}
\end{center}

\end{document}